\documentclass[10pt,a4paper]{amsart}
\usepackage{amsmath,amssymb,amsthm,graphicx}
% Retained minimal preamble for the Stacks Project chapter flat.tex (More on Flatness).
% Verbatim excerpt of the chapter preamble: Xy-pic setup (native preamble lines 14-18) and the native
% theorem declarations with the mathematical macros (native preamble lines 154-217). Copied mechanically;
% nothing added, removed or redefined. The wrapper class and the amsmath/amssymb/amsthm/graphicx base
% packages are supplied by the assembler, so the class/externaldocument part of the native preamble is omitted.
% For commutative diagrams we use Xy-pic
\usepackage[all]{xy}

% We use 2cell for 2-commutative diagrams.
\xyoption{2cell}
% Theorem environments.
%
\theoremstyle{plain}
\newtheorem{theorem}[subsection]{Theorem}
\newtheorem{proposition}[subsection]{Proposition}
\newtheorem{lemma}[subsection]{Lemma}

\theoremstyle{definition}
\newtheorem{definition}[subsection]{Definition}
\newtheorem{example}[subsection]{Example}
\newtheorem{exercise}[subsection]{Exercise}
\newtheorem{situation}[subsection]{Situation}

\theoremstyle{remark}
\newtheorem{remark}[subsection]{Remark}
\newtheorem{remarks}[subsection]{Remarks}

\numberwithin{equation}{subsection}

% Macros
%
\def\lim{\mathop{\mathrm{lim}}\nolimits}
\def\colim{\mathop{\mathrm{colim}}\nolimits}
\def\Spec{\mathop{\mathrm{Spec}}}
\def\Hom{\mathop{\mathrm{Hom}}\nolimits}
\def\Ext{\mathop{\mathrm{Ext}}\nolimits}
\def\SheafHom{\mathop{\mathcal{H}\!\mathit{om}}\nolimits}
\def\SheafExt{\mathop{\mathcal{E}\!\mathit{xt}}\nolimits}
\def\Sch{\mathit{Sch}}
\def\Mor{\mathop{\mathrm{Mor}}\nolimits}
\def\Ob{\mathop{\mathrm{Ob}}\nolimits}
\def\Sh{\mathop{\mathit{Sh}}\nolimits}
\def\NL{\mathop{N\!L}\nolimits}
\def\CH{\mathop{\mathrm{CH}}\nolimits}
\def\proetale{{pro\text{-}\acute{e}tale}}
\def\etale{{\acute{e}tale}}
\def\QCoh{\mathit{QCoh}}
\def\Ker{\mathop{\mathrm{Ker}}}
\def\Im{\mathop{\mathrm{Im}}}
\def\Coker{\mathop{\mathrm{Coker}}}
\def\Coim{\mathop{\mathrm{Coim}}}

% Boxtimes
%
\DeclareMathSymbol{\boxtimes}{\mathbin}{AMSa}{"02}

%
% Macros for moduli stacks/spaces
%
\def\QCohstack{\mathcal{QC}\!\mathit{oh}}
\def\Cohstack{\mathcal{C}\!\mathit{oh}}
\def\Spacesstack{\mathcal{S}\!\mathit{paces}}
\def\Quotfunctor{\mathrm{Quot}}
\def\Hilbfunctor{\mathrm{Hilb}}
\def\Curvesstack{\mathcal{C}\!\mathit{urves}}
\def\Polarizedstack{\mathcal{P}\!\mathit{olarized}}
\def\Complexesstack{\mathcal{C}\!\mathit{omplexes}}
% \Pic is the operator that assigns to X its picard group, usage \Pic(X)
% \Picardstack_{X/B} denotes the Picard stack of X over B
% \Picardfunctor_{X/B} denotes the Picard functor of X over B
\def\Pic{\mathop{\mathrm{Pic}}\nolimits}
\def\Picardstack{\mathcal{P}\!\mathit{ic}}
\def\Picardfunctor{\mathrm{Pic}}
\def\Deformationcategory{\mathcal{D}\!\mathit{ef}}

% --- appended 2026-10-12 (Chief): the project own definitions of the reference/slogan/history
% display environments, copied mechanically from the shared Stacks preamble lines 8-12
% (\usepackage{verbatim} plus the three \newenvironment lines, which the chapter preamble keeps
% verbatim from the project preamble). Nothing added, removed or redefined: these environments
% carry no mathematical content and the project itself typesets them as comments. Without these
% three lines a counted statement whose own text uses them fails with Environment slogan undefined.
\usepackage{verbatim}
\newenvironment{reference}{\comment}{\endcomment}
\newenvironment{slogan}{\comment}{\endcomment}
\newenvironment{history}{\comment}{\endcomment}
\usepackage{cleveref}
\crefname{originallemma}{original lemma}{original lemmas}
\Crefname{originallemma}{Original lemma}{Original lemmas}
\makeatletter
\newlabel{lemma-flat-after-blowing-up}{{081R}{source}{}{sourcepointer-1}{}}
\newlabel{lemma-flat-after-blowing-up@cref}{{[originallemma][0][]081R}{[1][0][]source}{}{}{}}
\makeatother
\crefname{originalproposition}{original proposition}{original propositions}
\Crefname{originalproposition}{Original proposition}{Original propositions}
\makeatletter
\newlabel{limits-proposition-approximate}{{01ZA}{source}{}{sourcepointer-2}{}}
\newlabel{limits-proposition-approximate@cref}{{[originalproposition][0][]01ZA}{[1][0][]source}{}{}{}}
\makeatother
\crefname{originallemma}{original lemma}{original lemmas}
\Crefname{originallemma}{Original lemma}{Original lemmas}
\makeatletter
\newlabel{limits-lemma-descend-finite-presentation}{{01ZM}{source}{}{sourcepointer-3}{}}
\newlabel{limits-lemma-descend-finite-presentation@cref}{{[originallemma][0][]01ZM}{[1][0][]source}{}{}{}}
\makeatother
\crefname{originallemma}{original lemma}{original lemmas}
\Crefname{originallemma}{Original lemma}{Original lemmas}
\makeatletter
\newlabel{limits-lemma-descend-opens}{{01Z4}{source}{}{sourcepointer-4}{}}
\newlabel{limits-lemma-descend-opens@cref}{{[originallemma][0][]01Z4}{[1][0][]source}{}{}{}}
\makeatother
\crefname{originallemma}{original lemma}{original lemmas}
\Crefname{originallemma}{Original lemma}{Original lemmas}
\makeatletter
\newlabel{limits-lemma-descend-flat-finite-presentation}{{04AI}{source}{}{sourcepointer-5}{}}
\newlabel{limits-lemma-descend-flat-finite-presentation@cref}{{[originallemma][0][]04AI}{[1][0][]source}{}{}{}}
\makeatother
\crefname{originallemma}{original lemma}{original lemmas}
\Crefname{originallemma}{Original lemma}{Original lemmas}
\makeatletter
\newlabel{morphisms-lemma-base-change-proper}{{01W4}{source}{}{sourcepointer-6}{}}
\newlabel{morphisms-lemma-base-change-proper@cref}{{[originallemma][0][]01W4}{[1][0][]source}{}{}{}}
\makeatother
\crefname{originallemma}{original lemma}{original lemmas}
\Crefname{originallemma}{Original lemma}{Original lemmas}
\makeatletter
\newlabel{morphisms-lemma-base-change-finite-presentation}{{01TS}{source}{}{sourcepointer-7}{}}
\newlabel{morphisms-lemma-base-change-finite-presentation@cref}{{[originallemma][0][]01TS}{[1][0][]source}{}{}{}}
\makeatother
\crefname{originallemma}{original lemma}{original lemmas}
\Crefname{originallemma}{Original lemma}{Original lemmas}
\makeatletter
\newlabel{morphisms-lemma-base-change-closed-immersion}{{01QR}{source}{}{sourcepointer-8}{}}
\newlabel{morphisms-lemma-base-change-closed-immersion@cref}{{[originallemma][0][]01QR}{[1][0][]source}{}{}{}}
\makeatother
\crefname{originallemma}{original lemma}{original lemmas}
\Crefname{originallemma}{Original lemma}{Original lemmas}
\makeatletter
\newlabel{morphisms-lemma-base-change-flat}{{01U9}{source}{}{sourcepointer-9}{}}
\newlabel{morphisms-lemma-base-change-flat@cref}{{[originallemma][0][]01U9}{[1][0][]source}{}{}{}}
\makeatother
\crefname{originallemma}{original lemma}{original lemmas}
\Crefname{originallemma}{Original lemma}{Original lemmas}
\makeatletter
\newlabel{lemma-base-change-almost-blow-up}{{0EV7}{source}{}{sourcepointer-10}{}}
\newlabel{lemma-base-change-almost-blow-up@cref}{{[originallemma][0][]0EV7}{[1][0][]source}{}{}{}}
\makeatother
\title{A finite presentation morphism that is flat over a quasi-compact open becomes flat after an almost blowup square}
\author{The Stacks Project Authors}
\date{stacks-project:flat.tex:local-snapshot-20261011; GNU FDL 1.2 or later}
\begin{document}
\maketitle
\noindent\textit{Editorial note.} All delivered mathematical text is the author's own, copied from the pinned Stacks Project chapter file flat.tex; no statement or proof step was rewritten, reordered or removed.\par\medskip
\noindent\textit{Editorial note.} The preamble is the chapter's own retained minimal preamble (sources/stacks-flat/retained-preamble.tex), a verbatim excerpt of the native Xy-pic setup and of the native theorem declarations and mathematical macros, followed by the project own definitions of the reference, slogan and history display environments (verbatim package plus three newenvironment lines) that the retained file carries at its end and that the project itself typesets as comments. No macro was invented or redefined, and the wrapper class and the amsmath/amssymb/amsthm/graphicx base packages are supplied by the assembler.\par\medskip
\noindent\textit{Editorial note.} The counted statement and proof contain no cross reference, no citation and no included file, so the delivered body is self contained.\par\medskip
\noindent\textit{Editorial note.} Canonical location of the counted proposition: Stacks tag 0EVC (sources/stacks-flat/pinned-original-source.tex line 11162).\par\medskip
\noindent\textit{Allowed standard prerequisites (editorial).}\begin{itemize}
\item The Noetherian case with blowup squares. Same chapter, Lemma labelled lemma-flat-after-blowing-up, cited at proof lines 11181-11185. Solves the Noetherian case and supplies the shape of the diagram.
\item Absolute Noetherian approximation and descent of finite presentation, opens and flatness. Limits, Proposition labelled limits-proposition-approximate, cited at proof lines 11189-11191, and Lemmas labelled limits-lemma-descend-finite-presentation, limits-lemma-descend-opens and limits-lemma-descend-flat-finite-presentation, cited at proof lines 11193-11201. Descends X -> Y, V and the flatness hypothesis to a Noetherian approximation.
\item Base change stability of almost blow up squares and of the relevant properties of morphisms. Same chapter, Lemma labelled lemma-base-change-almost-blow-up, cited at proof lines 11211-11213; Morphisms, Lemmas labelled morphisms-lemma-base-change-proper, morphisms-lemma-base-change-finite-presentation, morphisms-lemma-base-change-closed-immersion and morphisms-lemma-base-change-flat, cited at proof lines 11205-11211. Transports the Noetherian diagram to the given Y.
\end{itemize}\medskip
\ifdefined\hypertarget\hypertarget{sourcepointer-1}{}\fi
\noindent\textit{Source reference.} Original lemma 081R: sources/stacks-flat/pinned-original-source.tex, source lines 9123--9123. Same-chapter pointer into the author environment lemma at source line 9123; invoked at proof line 11183 of the counted proof. It is retained only as a navigation pointer; every load-bearing step is inside the delivered span.\par\medskip
\ifdefined\hypertarget\hypertarget{sourcepointer-2}{}\fi
\noindent\textit{Source reference.} Original proposition 01ZA: sources/stacks-chow/tags.txt, source lines 2585--2585. Cross-chapter standard prerequisite located in the project tag list staged in the local snapshot; invoked at proof line 11192 of the counted proof. It is retained only as a navigation pointer; every load-bearing step is inside the delivered span.\par\medskip
\ifdefined\hypertarget\hypertarget{sourcepointer-3}{}\fi
\noindent\textit{Source reference.} Original lemma 01ZM: sources/stacks-chow/tags.txt, source lines 2599--2599. Cross-chapter standard prerequisite located in the project tag list staged in the local snapshot; invoked at proof line 11195 of the counted proof. It is retained only as a navigation pointer; every load-bearing step is inside the delivered span.\par\medskip
\ifdefined\hypertarget\hypertarget{sourcepointer-4}{}\fi
\noindent\textit{Source reference.} Original lemma 01Z4: sources/stacks-chow/tags.txt, source lines 2579--2579. Cross-chapter standard prerequisite located in the project tag list staged in the local snapshot; invoked at proof line 11198 of the counted proof. It is retained only as a navigation pointer; every load-bearing step is inside the delivered span.\par\medskip
\ifdefined\hypertarget\hypertarget{sourcepointer-5}{}\fi
\noindent\textit{Source reference.} Original lemma 04AI: sources/stacks-chow/tags.txt, source lines 5617--5617. Cross-chapter standard prerequisite located in the project tag list staged in the local snapshot; invoked at proof line 11201 of the counted proof. It is retained only as a navigation pointer; every load-bearing step is inside the delivered span.\par\medskip
\ifdefined\hypertarget\hypertarget{sourcepointer-6}{}\fi
\noindent\textit{Source reference.} Original lemma 01W4: sources/stacks-chow/tags.txt, source lines 2469--2469. Cross-chapter standard prerequisite located in the project tag list staged in the local snapshot; invoked at proof line 11206 of the counted proof. It is retained only as a navigation pointer; every load-bearing step is inside the delivered span.\par\medskip
\ifdefined\hypertarget\hypertarget{sourcepointer-7}{}\fi
\noindent\textit{Source reference.} Original lemma 01TS: sources/stacks-chow/tags.txt, source lines 2385--2385. Cross-chapter standard prerequisite located in the project tag list staged in the local snapshot; invoked at proof line 11207 of the counted proof. It is retained only as a navigation pointer; every load-bearing step is inside the delivered span.\par\medskip
\ifdefined\hypertarget\hypertarget{sourcepointer-8}{}\fi
\noindent\textit{Source reference.} Original lemma 01QR: sources/stacks-chow/tags.txt, source lines 2276--2276. Cross-chapter standard prerequisite located in the project tag list staged in the local snapshot; invoked at proof line 11208 of the counted proof. It is retained only as a navigation pointer; every load-bearing step is inside the delivered span.\par\medskip
\ifdefined\hypertarget\hypertarget{sourcepointer-9}{}\fi
\noindent\textit{Source reference.} Original lemma 01U9: sources/stacks-chow/tags.txt, source lines 2402--2402. Cross-chapter standard prerequisite located in the project tag list staged in the local snapshot; invoked at proof line 11209 of the counted proof. It is retained only as a navigation pointer; every load-bearing step is inside the delivered span.\par\medskip
\ifdefined\hypertarget\hypertarget{sourcepointer-10}{}\fi
\noindent\textit{Source reference.} Original lemma 0EV7: sources/stacks-flat/pinned-original-source.tex, source lines 10987--10987. Same-chapter pointer into the author environment lemma at source line 10987; invoked at proof line 11212 of the counted proof. It is retained only as a navigation pointer; every load-bearing step is inside the delivered span.\par\medskip
\begin{lemma}
\label{lemma-flat-after-almost-blowing-up}
Let $Y$ be a quasi-compact and quasi-separated scheme.
Let $X$ be a scheme of finite presentation over $Y$.
Let $V \subset Y$ be a quasi-compact open such that
$X_V \to V$ is flat. Then there exist a commutative diagram
$$
\xymatrix{
E \ar[ddd] \ar[rd] & & & D \ar[lll] \ar[ddd] \ar[ld] \\
& Y' \ar[d] & X' \ar[l] \ar[d] \\
& Y & X \ar[l] \\
Z \ar[ru] & & & T \ar[lll] \ar[lu]
}
$$
whose right and left hand squares are almost blow up squares,
whose lower and top squares are cartesian, such that
$Z \cap V = \emptyset$, and
such that $X' \to Y'$ is flat (and of finite presentation).
\end{lemma}

\begin{proof}
If $Y$ is a Noetherian scheme, then this lemma follows immediately
from Lemma \ref{lemma-flat-after-blowing-up}
because in this case blow up squares are almost blow up squares
(we also use that strict transforms are blow ups).
The general case is reduced to the Noetherian case by
absolute Noetherian approximation.

\medskip\noindent
We may write $Y = \lim Y_i$ as a directed limit of an inverse system
of Noetherian schemes with affine transition morphisms, see
Limits, Proposition \ref{limits-proposition-approximate}.
We can find an index $i$ and a morphism $X_i \to Y_i$
of finite presentation whose base change to $Y$ is $X \to Y$.
See Limits, Lemmas \ref{limits-lemma-descend-finite-presentation}.
After increasing $i$ we may assume $V$ is the inverse image of
an open subscheme $V_i \subset Y_i$, see
Limits, Lemma \ref{limits-lemma-descend-opens}.
Finally, after increasing $i$ we may assume that $X_{i, V_i} \to V_i$
is flat, see
Limits, Lemma \ref{limits-lemma-descend-flat-finite-presentation}.
By the Noetherian case, we may construct a diagram as in the
lemma for $X_i \to Y_i \supset V_i$. The base change of this
diagram by $Y \to Y_i$ provides the solution. Use that base change
preserves properties of morphisms, see Morphisms, Lemmas
\ref{morphisms-lemma-base-change-proper},
\ref{morphisms-lemma-base-change-finite-presentation},
\ref{morphisms-lemma-base-change-closed-immersion}, and
\ref{morphisms-lemma-base-change-flat}
and that base change of an almost
blow up square is an almost blow up square, see
Lemma \ref{lemma-base-change-almost-blow-up}.
\end{proof}

\end{document}
